% Separate CP32 candidate. No global generation or connectedness is asserted.
\subsection{A nonempty full-inertia wall in every odd principal-cube system}
\label{subsec:ppav-all-odd-full-root-inertia}

Work first over an algebraically closed field \(k\) of characteristic two.
Let \(B\) be a nonordinary principally polarized abelian surface, choose
a symmetric representative \(\Theta_B\), and put
\(M=\mathcal O_B(3\Theta_B)\).  Let \(\pi:A\to B\) be an etale
isogeny of odd degree \(n\), put \(L=\pi^*M\), and set
\[
 W=H^0(A,L),\qquad W'=H^0(A^{(2)},L^{(2)}),\qquad N=9n.
 \tag{OW1}
\]
Relative Cartier and invariant volume forms define
\(\beta:\Lambda^2W\to W'\).  Alternation holds because the top
Cartier operator on a nonordinary abelian surface is zero.  Write
\(M_s=\beta(s,-)\) and \(B_\lambda(s,t)=\lambda(\beta(s,t))\).
We use the separately proved algebraic statement
Theorem~\ref{thm:DE-Moore-Fermat-nondivisibility}: for
\[
 D_c(x)=\det\bigl(c_{i-j}x_{i+j}\bigr)_{i,j\in\mathbf F_3^2},
 \qquad \mathcal F(x)=\sum_{i\in\mathbf F_3^2}x_i^3,
 \tag{OW2}
\]
the divisibility \(\mathcal F\mid D_c\) holds if and only if
\(D_c\) is identically zero.  This is an exact universal polynomial
statement, not an inference from values at finitely many points.

\begin{lemma}[One Fermat coordinate system for both kinds of twisted blocks]
\label{lem:odd-cover-common-Fermat-coordinates}
\status{Candidate; complete proof supplied for independent audit.}
Fix the symmetric level-three theta basis \((e_i)_{i\in\mathbf F_3^2}\)
on \(H^0(B,M)\) from
Lemma~\ref{lem:ppav-symmetric-level3-addition-gauge}, and its actual
Frobenius-twisted basis \((e_i')\) on \(H^0(B^{(2)},M^{(2)})\).
Use coordinates \(s_i\) in the first basis and \(x_i=\lambda(e_i')\)
on the dual of the second.  Then the central multiplication matrix
is symmetric alternating, with both kernel coordinate vectors equal
to \(s\) on its rank-eight open.  Its relative wall is
\(\mathcal F(s)=0\).

For each nontrivial odd-order \(P\in\operatorname{Pic}^0(B)(k)\),
the following two determinant polynomials, in these same fixed
coordinates \(s_i\), are nonzero and not divisible by \(\mathcal F(s)\):
the cross block for \(MP,MP^{-1}\), evaluated at
\(\lambda(e_i')=s_i\); and the multiplication map
\[
 H^0(B,MP)\longrightarrow H^0(B^{(2)},M^{(2)}Q^{(2)}),
 \qquad t\longmapsto C_{\rm rel}(st\omega_B)/\omega_B',
 \qquad Q^2=P
 \tag{OW3}
\]
where \(Q\) is the unique half of \(P\) in its odd cyclic subgroup.
\end{lemma}
\begin{proof}
The symmetric splitting gives the standard finite Heisenberg action
in precisely the displayed basis.  The proof of
Theorem~\ref{thm:ppav-level3-Cartier-symmetric} therefore applies
without changing that basis: its equivariant identification sends
\(e_i'\) to \(e_i^*\), and its matrix is the identity.  The resulting
trilinear tensor is fully symmetric.  Since it is alternating in
two arguments, the central multiplication matrix is alternating and
has right kernel \(ks\), hence left kernel with coordinate vector
\(s\), whenever its rank is eight.  Pairing that left kernel with
\(s^{[2]}\) gives exactly \(\sum_i s_i^3\).

For the cross block choose \(a\) with \(\phi_M(a)=P\).  The exact
identity (TCM5), with inputs \(e_i(z+a),e_j(z-a)\) and output basis
\(e_k'(z')\), has coefficients
\[
 c\,c_{(i-j)/2}(a)e_{(i+j)/2}'(z'),
 \qquad
 c_h(a)=\mathcal T_{a_0}(e_h^2)(a),
 \tag{OW4}
\]
after one common fibre trivialization; \(c\ne0\) is a common scalar.
Here \(a_0\in B[2](k)\) is fixed, as in
Proposition~\ref{prop:ppav-Cartier-cross-Moore-identification}.
Reindexing \(i=2u,j=2v\) makes the matrix of the contraction by
\(\lambda\) exactly \(c\,c_{u-v}(a)x_{u+v}\).  Setting \(x_i=s_i\)
therefore gives a scalar multiple of \(D_{c(a)}(s)\).
Theorem~\ref{thm:ppav-non-two-torsion-cross-Cartier} makes this
determinant nonzero because \(P^2\ne\mathcal O_B\).  The universal
nondivisibility theorem then excludes \(\mathcal F(s)\) as a factor.

For (OW3), choose \(u\) with \(\phi_M(u)=Q\).  Retain the original
first input basis \(e_i(z)\).  Choose for the second input and output
the bases obtained from \(e_j(z+2u)\) and
\(e_k'(z'+F_B(u))\), respectively.  These have the required line
bundles because
\[
 t_{2u}^*M\simeq MP,\qquad
 t_{F_B(u)}^*M^{(2)}\simeq M^{(2)}Q^{(2)},\qquad
 F_B^*(M^{(2)}Q^{(2)})=M^2P.
 \tag{OW5}
\]
Put \(w=z+u\) and apply (TCM5) at parameter \(a=-u\).
The two balanced inputs are then exactly \(e_i(z)\) and
\(e_j(z+2u)\).  Naturality of relative Cartier under translation
gives
\[
 C_{\rm rel}\bigl(e_i(z)e_j(z+2u)\omega_B\bigr)
 =c\,c_{(i-j)/2}(-u)\,
       e_{(i+j)/2}'(z'+F_B(u))\omega_B'.
 \tag{OW6}
\]
Thus multiplication by the original section \(\sum_i s_i e_i\)
has output-row \(k\), input-column \(j\) entry
\[
 c\,c_{k-j}(-u)s_{2k-j}
     =c\,c_{k-j}(-u)s_{-(k+j)}.
 \tag{OW7}
\]
This is the same Moore matrix with the sole variable change
\(x_v=s_{-v}\), an index permutation preserving \(\mathcal F\)
exactly.  The parameter corresponds to \(Q^{-1}\), which is again
nontrivial of odd order, so its determinant is nonzero by the same
cross-block theorem.  Universal nondivisibility proves the claim.

In particular, translating to opposite twists has not introduced an
arbitrary linear change in the \(s_i\): the original first input in
(OW6) is literally \(e_i(z)\).  The two-rank-one correction \(a_0\)
changes only the coefficient vector \(c(-u)\), not these variables.
All square-isomorphism and line-bundle identifications above are the
normalized ones in (TCM5); at the fixed parameter their ambiguity is
one common nonzero scalar.  No independent, parameter-dependent row
gauges or unproved invariance of a Fermat polynomial are used.
\end{proof}

\begin{proposition}[The maximal common-rank open meets the relative wall]
\label{prop:odd-cover-common-rank-wall-nonempty}
\status{Candidate; complete proof supplied for independent audit.}
On the maximal-rank locus for \(M_s\), let \(k\lambda_s\) be its left
kernel, and define the open subset
\[
 U_{\max}=\{[s]\in\mathbf P(W):
      \operatorname{rank}M_s=N-1,\
      \operatorname{rank}B_{\lambda_s}=N-1\}.
 \tag{OW8}
\]
Then
\[
 H_{\max}=\{[s]\in U_{\max}:\lambda_s(s^{[2]})=0\}
 \tag{OW9}
\]
is nonempty.  This assertion concerns the maximal common-rank open,
not the smaller smooth and Frobenius-bijective open used to prove
generic sharpness in Proposition~\ref{prop:DE-common-Cartier-incidence}.
\end{proposition}
\begin{proof}
The odd-character decomposition gives
\[
 W=\bigoplus_\chi H^0(B,MP_\chi),\qquad
 W_0=H^0(B,M),\qquad P_1=\mathcal O_B.
 \tag{OW10}
\]
Restrict \(s\) to \(W_0\).  Multiplication by such a section sends
the \(P_\chi\) block to the target \(Q_\chi^{(2)}\) block, where
\(Q_\chi^2=P_\chi\).  Thus its nontrivial blocks are the maps (OW3),
with the target characters permuted by their unique halves.

The central four-parameter Pfaffian cubic \(q\) is nonzero by
Proposition~\ref{prop:DE-A3-ppav-degree-one}.  Lemma
\ref{lem:nonordinary-uniform-Q-stable-loss-candidate} proves that
\(q\) is not proportional to \(\mathcal F\), and that the central
matrix has rank eight on the dense open \(\mathcal F\setminus\{q=0\}\).
The projective Fermat cubic is smooth, since its partial derivatives
are \(s_i^2\), and is geometrically integral.

Every nontrivial multiplication determinant is nonzero on this
integral cubic by the preceding lemma.  Intersect their finitely
many nonvanishing opens with \(q\ne0\).  On that intersection the
total multiplication matrix has rank \(8+9(n-1)=N-1\), and its
left kernel is supported on the trivial target block with coordinates
\(\lambda_s(e_i')=s_i\).

For this same left kernel, the dual alternating form has the central
rank-eight block and the nontrivial inverse-character cross blocks.
The preceding lemma makes each cross determinant nonzero on the same
Fermat cubic, in the same \(s_i\).  Intersect those finitely many
nonempty opens as well.  The resulting open in the integral cubic
is nonempty.  Its points have
\(\operatorname{rank}B_{\lambda_s}=8+18(n-1)/2=N-1\).
They consequently belong to \(U_{\max}\), and their wall equation is
exactly \(\lambda_s(s^{[2]})=\sum_i s_i^3=0\).  This proves (OW9).
No smoothness of the divisor defined by such a section is required.
\end{proof}

\begin{theorem}[Full hyperplane inertia for every odd cubic pullback]
\label{thm:odd-cover-full-hyperplane-inertia}
\status{Candidate; complete proof supplied for independent audit.}
For the family (OW1), the geometric generic relative Lang image
contains a full hyperplane radical of order \(2^{N-2}\).
More precisely, every component of \(H_{\max}\) is smooth and
reduced in \(U_{\max}\).  At its geometric generic point the
relative coherent Frobenius has first and stable rank \(N-2\)
and determinant order one; its divisorial inertia contains
\[
 \operatorname{Hom}(\mathbf F_2,\mathbf F_2^{N-2})
 \subset\operatorname{GL}_{N-1}(\mathbf F_2).
 \tag{OW11}
\]
The same containment holds for the image of the relative local system
on the dense smooth-member, maximal-rank open.
\end{theorem}
\begin{proof}
We check the coherent model on the entire coefficient projective
space, including its nonsmooth members.  The universal divisor is
defined by the universal nonzero section of
\(L\boxtimes\mathcal O_{\mathbf P(W)}(1)\).  Its hypersurface
sequence and \(H^1(A,L^{-1})=0\) identify the quotient of relative
\(H^1(\mathcal O)\) by the constant \(H^1(A,\mathcal O_A)\) with
the kernel
\[
 \mathcal E=
 \ker\bigl(W^*\otimes\mathcal O_{\mathbf P(W)}(-1)
                      \longrightarrow\mathcal O_{\mathbf P(W)}\bigr)
 \simeq\Omega^1_{\mathbf P(W)}.
 \tag{OW12}
\]
The last map is evaluation at \(s\) and is everywhere surjective.
Here Serre duality uses the fixed invariant volume form.  Coherent
Frobenius is regular on this bundle over the whole coefficient space.
The same multiplication/Cartier duality as in
Proposition~\ref{prop:birational-Cartier-kernel-wall-inertia}
gives on a coefficient chart its actual formula
\[
 F_s=M_s^t:(s^{[2]})^\perp\longrightarrow s^\perp.
 \tag{OW13}
\]
This construction is a hypersurface-cohomology calculation; it
requires a nonzero section, but not smoothness of its zero divisor.

All three hypotheses of that proposition now hold on \(U_{\max}\).
The preceding proposition supplies its missing nonemptiness
hypothesis.  It follows that \(H_{\max}\) is smooth and reduced,
that the Euler determinant has order one there, and that its
linearized kernel map has differential rank \(N-2\).  The squared-frame
lemma gives stable loss exactly one.  Finally the full
strict-henselian Artin--Schreier argument of
Lemma~\ref{lem:moving-kernel-full-root-inertia-candidate} gives
(OW11), including the possible inseparability of the residue extension.

For the last sentence, the Frobenius-bijective locus of \(\mathcal E\)
is a nonempty open by
Proposition~\ref{prop:DE-common-Cartier-incidence}; that proposition
also supplies smooth sharp members.  The generic relative Lang
local system is finite etale on that open.  Restricting it to its
dense smooth-member open leaves its image unchanged.  Indeed, over
a normal integral base every connected finite etale cover is
normal and integral, and remains connected after restriction to a
nonempty open; equivalently the map of fundamental groups from
the smaller open is surjective.  Both local systems have the same
generic representation.  Thus the divisorial inertia just computed
lies in the smooth-family image even if the selected generic wall
divisor parametrizes nonsmooth curves.  No specialization isomorphism
between stable fibres across a rank drop is asserted.
\end{proof}

\begin{corollary}[Every fixed odd principal-cube degree]
\label{cor:fixed-ppav-odd-full-hyperplane-inertia}
\status{Candidate; complete proof supplied for independent audit.}
Let \((A,\Theta)\) be a fixed nonordinary principally polarized
abelian surface over a perfect field of characteristic two.
For every odd integer \(d\ge1\), the geometric relative Lang image
for the complete system \(|3d\Theta|\) contains a full hyperplane
radical of order \(2^{9d^2-2}\) on its rank-\((9d^2-1)\) space.
\end{corollary}
\begin{proof}
After geometric base change, the exact maximal-isotropic quotient
and symmetric descent in
Corollary~\ref{cor:fixed-ppav-all-odd-maximal-dual-Cartier}
write \(\mathcal O_A(3d\Theta)=\pi^*(M_0^3)\), where \(\pi\)
is etale of degree \(d^2\) and \(M_0\) is principally polarizing
on a nonordinary surface.  Apply the theorem with \(n=d^2\).
The case \(d=1\) is included and agrees with the earlier
degree-one full-inertia theorem.  The conclusion is geometric;
any arithmetic image containing it inherits this subgroup.
\end{proof}

\begin{remark}[Scope and the remaining global question]
\label{rem:odd-full-hyperplane-inertia-scope}
\status{Candidate scope; no new full-group degree claimed.}
This proves that the common Cartier wall is nonempty in every
stated fixed odd system and that its inertia is a complete
hyperplane radical.  It does not prove connectedness of the
nonzero relative Lang cover, irreducibility of the global module,
or normal generation by these particular radicals.  Normal
generation by all divisorial inertia is a different assertion;
the other divisors have not been identified with this wall.
The previously retained parabolic obstruction is not removed by
the mere existence of one such subgroup.  No new fixed-surface
full-image degree is asserted here.

The new coordinate and geometric comparisons are proved directly.
The only finite algebraic input is the separately retained universal
Moore Fermat-nondivisibility certificate, whose proof and computation
scope remain those of its own subsection.  No new numerical run
is needed or claimed for the present deductions.
\end{remark}
